\documentclass[11pt]{article}
\usepackage{amsmath,amssymb,amsthm}
\usepackage[margin=1in]{geometry}

\title{Audited Communication Reduction for Resolution}
\author{Inacio F. Vasquez}
\date{}

\newtheorem{theorem}{Theorem}
\newtheorem{corollary}{Corollary}
\newtheorem{boundary}{Boundary}

\begin{document}
\maketitle

\section*{Audit correction}

Let $F$ be an unsatisfiable CNF and let $\mathrm{Search}_F$ denote its
falsified-clause search problem.  Fix any partition of the variables between
two communicating parties.

The previously stated implication
\[
IC_\mu(\mathrm{Search}_F)\le O(\log L(F))
\]
was not justified.  Its proof sketch implicitly balanced an arbitrary
resolution DAG to logarithmic depth.  General resolution proofs do not admit
such a size-preserving balancing theorem.

The standard proof-complexity correspondence is instead DAG-preserving:
a size-$s$ resolution refutation corresponds to a size-$s$ conjunction-DAG
solving $\mathrm{Search}_F$, and feasible interpolation gives a monotone
circuit of size $O(s)$ for the associated unsatisfiability certificate.
Ordinary deterministic communication protocols correspond to formula/tree
objects, not arbitrary circuit/DAG size.

\section*{Valid tree/depth transfer}

\begin{theorem}
If $F$ has a tree-like resolution refutation of depth $d$, then for every
variable partition there is a deterministic two-party protocol solving
$\mathrm{Search}_F$ with communication cost at most $d$.
\end{theorem}

\begin{proof}
A depth-$d$ tree-like resolution refutation is equivalent to a depth-$d$
decision tree solving $\mathrm{Search}_F$.  The two parties simulate that
decision tree.  At each queried variable, the party holding that variable
communicates its value, using at most one bit per level.  Thus the total
communication is at most $d$.
\end{proof}

\begin{corollary}
For every input distribution $\mu$,
\[
IC_\mu(\mathrm{Search}_F)\le d
\]
whenever $F$ has a tree-like resolution refutation of depth $d$.
\end{corollary}

\section*{DAG-like size boundary}

\begin{boundary}
For general DAG-like resolution, a lower bound on ordinary transcript
information complexity does not by itself imply
\[
L(F)\ge 2^{\Omega(IC_\mu(\mathrm{Search}_F))}.
\]
To obtain a general resolution-size lower bound through the communication
route, one needs a lower bound for a DAG-compatible object---for example the
associated monotone circuit / rectangle-DAG model, or a proved information
measure that lower-bounds that DAG-like model.  No such transfer theorem is
claimed here.
\end{boundary}

\section*{Next missing object}

The next proof obligation is to identify and prove a DAG-compatible
information-complexity lower bound for the Tseitin search family, or else to
replace this route by an existing lifting theorem that transfers an
appropriate query/width lower bound to monotone-circuit and resolution size.

\end{document}
